\documentclass[reqno,11pt]{amsart}

\usepackage[margin=1in]{geometry}
\usepackage[T1]{fontenc}
\usepackage{lmodern}
\usepackage{microtype}
\usepackage{amsmath,amssymb,amsthm,mathtools}
\usepackage[dvipsnames]{xcolor}
\usepackage[most]{tcolorbox}
\usepackage{tikz}
\usepackage{enumitem}
\usepackage[hidelinks]{hyperref}
\hypersetup{
  pdftitle={Lecture 4, MATH 327: Subgroups and Greatest Common Divisors},
  pdfauthor={Manish Patnaik},
  pdfsubject={MATH 327 course notes}
}

\definecolor{courseblue}{HTML}{1F4E79}
\definecolor{lightblue}{HTML}{EAF2F8}
\definecolor{darkgray}{HTML}{333333}
\definecolor{definitiongreen}{HTML}{EEF8EE}
\definecolor{definitionborder}{HTML}{8FB58F}
\definecolor{lightpink}{HTML}{FCE8EF}

\setlength{\parindent}{0pt}
\setlength{\parskip}{0.65em}
\setlist[itemize]{topsep=0.25em,itemsep=0.15em}

% Dotted leaders in the table of contents.
\makeatletter
\def\@tocline#1#2#3#4#5#6#7{\relax
  \ifnum #1>\c@tocdepth
  \else
    \par\addpenalty\@secpenalty\addvspace{#2}%
    \begingroup\hyphenpenalty\@M
    \@ifempty{#4}{%
      \@tempdima\csname r@tocindent\number#1\endcsname\relax
    }{%
      \@tempdima#4\relax
    }%
    \parindent\z@\leftskip#3\relax\advance\leftskip\@tempdima\relax
    \rightskip\@pnumwidth plus4em\parfillskip-\@pnumwidth
    #5\leavevmode\hskip-\@tempdima
    \ifcase #1
      \or\or\hskip 1em\or\hskip 2em\else\hskip 3em
    \fi
    #6\nobreak\relax
    \dotfill\hbox to\@pnumwidth{\@tocpagenum{#7}}\par
    \nobreak
    \endgroup
  \fi}
\makeatother

\theoremstyle{definition}
\newtheorem{definition}{Definition}[section]
\tcolorboxenvironment{definition}{
  enhanced,
  breakable,
  colback=definitiongreen,
  colframe=definitionborder,
  boxrule=0.5pt,
  arc=1mm,
  left=2mm,
  right=2mm,
  top=1mm,
  bottom=1mm,
  before skip=0.8em,
  after skip=0.8em
}
\newtheorem{example}{Example}[section]
\theoremstyle{plain}
\newtheorem{theorem}{Theorem}[section]
\newtheorem{lemma}{Lemma}[section]
\newtheorem{proposition}{Proposition}[section]
\theoremstyle{remark}
\newtheorem*{remark}{Remark}
\newtheorem{exercise}[theorem]{Exercise}

\tcolorboxenvironment{exercise}{
	enhanced,
	colback=lightpink,
	colframe=WildStrawberry,
	boxrule=0.7pt,
	arc=2pt,
	left=6pt,
	right=6pt,
	top=5pt,
	bottom=5pt
}


\numberwithin{equation}{section}

\newcommand{\Z}{\mathbb{Z}}
\newcommand{\R}{\mathbb{R}}
\newcommand{\GL}{\operatorname{GL}}
\newcommand{\id}{\operatorname{id}}
\newcommand{\be}{\begin{eqnarray}}
\newcommand{\ee}{\end{eqnarray}}

\begin{document}

\begin{center}
  \colorbox{lightblue}{%
    \begin{minipage}{0.94\textwidth}
      \vspace{0.55em}
      \begin{tabular*}{\textwidth}{@{\extracolsep{\fill}}lr}
        {\color{courseblue}\bfseries\LARGE Lecture 4} &
        {\color{courseblue}\bfseries\LARGE MATH 327}\\[0.35em]
         & \textbf{Date: Sep. 11, 2026}\\
        & \textbf{Instructor: Manish Patnaik}
      \end{tabular*}
      \vspace{0.45em}
    \end{minipage}}
\end{center}

\setcounter{tocdepth}{1}
\tableofcontents

\section{Recollections on Subgroups}

\begin{definition}[Subgroup]
Let $(G,\cdot,e)$ be a group. A subset $H\subseteq G$ is a \emph{subgroup}
if it satisfies the following conditions:
\begin{enumerate}[label=\textbf{(\arabic*)}]
\item $e\in H$.
\item $H$ is closed under products: if $h_1,h_2\in H$, then $h_1h_2\in H$.
\item $H$ is closed under inverses: if $h\in H$, then $h^{-1}\in H$,
where the inverse is taken in $G$.
\end{enumerate}
\end{definition}

A subgroup is itself a group under the same operation, with the same
identity as $G$. Associativity is inherited from $G$. Recall that we also proved the following result earlier in Lecture 2.

\begin{proposition}[Finite subgroup criterion]
If $H\subseteq G$ is finite, contains $e$, and is closed under products,
then $H$ is a subgroup. In this case, closure under inverses follows from
the other two conditions.
\end{proposition}


\begin{remark}
For an infinite subset, closure under inverses must still be checked.
For example, $\{0,1,2,\ldots\}\subseteq(\Z,+,0)$ contains the identity
and is closed under addition, but it is not closed under additive inverses.
\end{remark}

\section{Generated subgroups}

Recall that we have also already introduced the following notion.

\begin{definition}[Cyclic subgroup]
Let $G$ be a group, with identity denoted by $1$, and let $x\in G$.
The \emph{cyclic subgroup generated by $x$} is
\be
\langle x\rangle=\{x^m:m\in\Z\}
 =\{1,x^{\pm1},x^{\pm2},x^{\pm3},\ldots\}.
\ee
\end{definition}

\newpage
We generalize this definition as follows.

\begin{definition}[Subgroup generated by a subset]
For a subset $W\subseteq G$, let $\langle W\rangle$ be the set of all
finite products
\be
w_1w_2\cdots w_k,\qquad k\geq0,
\ee
where each $w_i$ belongs to $W$ or has its inverse in $W$.
When $k=0$, the product is the identity $1$. This is called the
\emph{subgroup generated by $W$}.
\end{definition}

Indeed, concatenating two such products gives another such product. As for inverses, let's mention the following important but simple identity that you may remember from a course in linear algebra.
\be
(w_1w_2\cdots w_k)^{-1}=w_k^{-1}\cdots w_2^{-1}w_1^{-1}
\ee
The product on the right is again of the required form. Hence
$\langle W\rangle$ is a subgroup.

\begin{remark}
The subgroup $\langle W\rangle$ is the \emph{smallest subgroup containing
$W$}. Indeed, it contains $W$ by taking products of length one. If $H$ is
any subgroup containing $W$, then closure under inverses and products
implies that $H$ contains every product used to define $\langle W\rangle$.
Thus $\langle W\rangle\subseteq H$.
\end{remark}

\begin{example}[Subgroups of the integers]
In the additive group $(\Z,+,0)$, the even integers, $\{0\}$, and $\Z$
are all subgroups. More generally, for $a\in\Z$,
\be
\Z a=\langle a\rangle=\{ma:m\in\Z\}
 =\{0,\pm a,\pm2a,\pm3a,\ldots\}
\ee
is the subgroup consisting of all multiples of $a$.
Here the operation is addition, so powers in multiplicative notation
become integer multiples in additive notation.
\end{example}

\begin{definition}
For any group $G$, the set $\{e\}$ is a subgroup, called the
\emph{trivial subgroup}. The whole group $G$ is also a subgroup of itself. A subgroup $H$ is called a \emph{proper subgroup} if $H\neq G$.
If $\{e\}\subsetneq H\subsetneq G$, then $H$ is called a \emph{nontrivial proper
subgroup}.

\end{definition}

Intersecting subgroups provides another useful way to manufacture new groups.


\begin{proposition}
If $H$ and $K$ are subgroups of $G$, then $H\cap K$ is a subgroup of $G$.
\end{proposition}

\begin{proof}
Since $e\in H$ and $e\in K$, we have $e\in H\cap K$.
If $a,b\in H\cap K$, then $ab\in H$ and $ab\in K$, because both
subgroups are closed under products. Thus $ab\in H\cap K$.
Similarly, if $a\in H\cap K$, then $a^{-1}\in H$ and $a^{-1}\in K$,
so $a^{-1}\in H\cap K$.
\end{proof}

\begin{remark}
The union $H\cup K$ need not be a subgroup. For example,
$2\Z\cup3\Z$ contains $2$ and $3$, but does not contain their sum $5$.
\end{remark}

\section{Subgroups of \texorpdfstring{$\Z$}{Z} and greatest common divisors}

\begin{lemma}
Every subgroup of $(\Z,+,0)$ is either $\{0\}$ or $\Z d$ for some
positive integer $d$.
\end{lemma}

\begin{proof}
Let $H\subseteq\Z$ be a subgroup with $H\neq\{0\}$. Since $H$ contains
a nonzero integer and its negative, it contains a positive integer.
By the well-ordering principle, $H$ has a smallest positive element $d$.
We claim that $H=\Z d$.

First, $\Z d\subseteq H$, since $d\in H$ and $H$ is closed under
addition and additive inverses.

For the reverse inclusion, let $m\in H$. By division with remainder,
write
\be
m=qd+r,\qquad q\in\Z,\quad 0\leq r<d.
\ee
Since $m,qd\in H$, we have $r=m-qd\in H$. If $r>0$, this contradicts
the choice of $d$ as the smallest positive element of $H$. Thus $r=0$,
so $m=qd\in\Z d$. Hence $H\subseteq\Z d$, proving the claim.
\end{proof}

\subsection*{Greatest common divisors}

Let $a,b\in\Z$. Consider the set of all integer linear combinations
of $a$ and $b$:
\be
\Z a+\Z b=\{ma+nb:m,n\in\Z\}.
\ee
This is a subgroup of $(\Z,+,0)$: it contains $0$, and sums and negatives
of integer linear combinations are again integer linear combinations.
Therefore the preceding lemma gives
\be
\Z a+\Z b=\Z d
\ee
for some $d\geq0$. If $a$ and $b$ are not both zero, then $d>0$.
We will show that this positive generator is $\gcd(a,b)$.

\begin{lemma}[Divisibility and inclusion]\label{lem:divisibility}
For integers $r,s$,
\be
r\mid s\quad\Longleftrightarrow\quad\Z s\subseteq\Z r.
\ee
\end{lemma}

\begin{proof}
If $s=kr$ for some $k\in\Z$, then every multiple of $s$ is a multiple
of $r$. Conversely, if $\Z s\subseteq\Z r$, then $s\in\Z r$, so $r\mid s$.
\end{proof}

\begin{proposition}
Let $a,b\in\Z$ be not both zero, and let $d>0$ satisfy
$\Z a+\Z b=\Z d$. Then $d=\gcd(a,b)$. In particular,
\be
\Z a+\Z b=\Z\gcd(a,b).
\ee
\end{proposition}

\begin{proof}
We check two properties:
\begin{enumerate}[label=\textbf{(\arabic*)}]
\item $d$ divides both $a$ and $b$.
\item Every common divisor of $a$ and $b$ divides $d$.
\end{enumerate}
For the first property, the inclusions
\be
\Z a\subseteq\Z a+\Z b=\Z d,
\qquad
\Z b\subseteq\Z a+\Z b=\Z d
\ee
imply $d\mid a$ and $d\mid b$ by Lemma~\ref{lem:divisibility}.

For the second property, suppose $\ell\mid a$ and $\ell\mid b$.
Then $\Z a\subseteq\Z\ell$ and $\Z b\subseteq\Z\ell$. Adding elements
from these two subgroups gives
\be
\Z d=\Z a+\Z b\subseteq\Z\ell.
\ee
By Lemma~\ref{lem:divisibility}, $\ell\mid d$.
Since $d>0$, every positive common divisor is therefore at most $d$.
Thus $d$ is the greatest common divisor of $a$ and $b$.
\end{proof}

A corollary of the above result is the following useful result.

\begin{proposition}[B\'ezout's identity]
Let $a,b\in\Z$ be not both zero. There exist integers $m,n$ such that
\be
\gcd(a,b)=ma+nb.
\ee
\end{proposition}

\begin{proof}
Set $d=\gcd(a,b)$. Since $\Z a+\Z b=\Z d$ and $d\in\Z d$, we have
$d\in\Z a+\Z b$. By the definition of this subgroup, $d=ma+nb$ for some
$m,n\in\Z$.
\end{proof}

\end{document}
